% =====================================================================
%  Control de versiones del documento, revisiones y aprobaciones
% =====================================================================
\clearpage
\section*{Control de versiones}
\phantomsection
\addcontentsline{toc}{part}{Control de versiones}

\begin{instruccion}
Registre aquí la evolución de \emph{este documento}, distinguiendo tres cosas: el
\textbf{historial} de cambios (quién modificó qué), la \textbf{revisión} (alguien leyó el
documento y emitió observaciones) y la \textbf{aprobación formal} (alguien con autoridad acepta
el documento como válido). Cada una tiene su propio cuadro.
\end{instruccion}

\begin{notanormativa}[Alcance de este apartado]
Este control de versiones cubre únicamente el \emph{documento}. El control de versiones de los
demás activos de prueba —software bajo prueba, casos, procedimientos, guiones, datos y
configuración del entorno— se define en el apartado~\ref{sec:configuracion} y se registra en
cada ejecución (\ref{anx:ejecuciones}).
\end{notanormativa}

\begin{observacion}[Revisión y aprobación no son lo mismo]
\obsproblema{La plantilla anterior resolvía las revisiones con un único cuadro cuya columna
«Estatus» mezclaba el resultado de una revisión con la aprobación formal del documento, y
limitaba todo el control de versiones al propio archivo Word.}
\obscambio{El apartado se divide en historial, revisiones y aprobación formal; el control de
versiones del resto de los activos pasa al apartado~\ref{sec:configuracion}.}
\obsjustificacion{Una revisión puede cerrarse con observaciones sin que el documento quede
aprobado, y un resultado de prueba no es reproducible si sólo se versiona el documento que lo
describe.}
\obsfuente{OBS §2 y §3}
\end{observacion}

% ------------------------------ Historial ----------------------------
\subsection*{Historial de cambios}

\begin{table}[H]
\centering
\caption{Historial de cambios del documento.}
\label{tab:historial}
\begin{tabularx}{\textwidth}{|C{2.5cm}|C{1.8cm}|L{3.4cm}|Y|}
  \hline
  \filaenc \textbf{Fecha} & \textbf{Versión} & \textbf{Autor} & \textbf{Descripción del cambio} \\ \hline
  \campo{dd/mm/aaaa} & \campo{0.1} & \campo{Nombre} & \campo{Qué se agregó, modificó o eliminó} \\ \hline
  & & & \\ \hline
  & & & \\ \hline
  & & & \\ \hline
\end{tabularx}
\end{table}

El Cuadro~\ref{tab:historial} deja constancia de cada versión emitida del plan; debe
actualizarse \emph{antes} de cada entrega, no después.

% ------------------------------ Revisiones ---------------------------
\subsection*{Revisiones}

\begin{table}[H]
\centering
\caption{Revisiones realizadas al documento.}
\label{tab:revisiones}
\begin{tabularx}{\textwidth}{|C{2.5cm}|C{1.5cm}|L{2.9cm}|Y|C{2.4cm}|}
  \hline
  \filaenc \textbf{Fecha} & \textbf{Versión} & \textbf{Revisor} & \textbf{Observaciones} & \textbf{Resultado} \\ \hline
  \campo{dd/mm/aaaa} & \campo{0.1} & \campo{Nombre y rol} & \campo{Hallazgos de la revisión} & \campo{Ver nota} \\ \hline
  & & & & \\ \hline
  & & & & \\ \hline
\end{tabularx}
\notatabla{Resultado: \textit{Sin observaciones} · \textit{Con observaciones menores} · \textit{Con observaciones mayores}. La revisión no equivale a la aprobación.}
\end{table}

El Cuadro~\ref{tab:revisiones} registra quién leyó el documento y qué observó. Una revisión
puede cerrarse sin que el documento quede aprobado.

% ----------------------------- Aprobaciones --------------------------
\subsection*{Aprobación formal}

\begin{table}[H]
\centering
\caption{Aprobación formal del documento.}
\label{tab:aprobaciones}
\begin{tabularx}{\textwidth}{|C{2.5cm}|C{1.5cm}|L{3.4cm}|Y|C{2.2cm}|}
  \hline
  \filaenc \textbf{Fecha} & \textbf{Versión} & \textbf{Autoridad que aprueba} & \textbf{Ámbito de la aprobación} & \textbf{Estatus} \\ \hline
  \campo{dd/mm/aaaa} & \campo{1.0} & \autoridadAprobacion & \campo{Plan completo / apartado} & \campo{Ver nota} \\ \hline
  & & & & \\ \hline
\end{tabularx}
\notatabla{Estatus: \textit{Aprobado} · \textit{Aprobado con correcciones} · \textit{No aprobado}. Corresponde a la tarea TP8 «Lograr consenso sobre el plan» de \normap{2}.}
\end{table}

El Cuadro~\ref{tab:aprobaciones} identifica a quién autorizó el plan y sobre qué versión lo
hizo; sin esa firma, el plan permanece en borrador y no debería dirigir la ejecución.

\begin{politicaacademica}
Durante el curso, la autoridad de aprobación del plan y de cualquier desviación respecto de la
estrategia acordada es \autoridadAprobacion.\ifversionfinal\else\ Al reutilizar la plantilla
fuera del curso, sustituya ese valor en \texttt{config/metadata.tex} por el rol que corresponda
en la organización (líder de pruebas, responsable de calidad, patrocinador del proyecto).\fi
\end{politicaacademica}
